% Generated draft: requires mathematical and source review.
\documentclass[11pt]{article}
\usepackage[margin=25mm]{geometry}
\usepackage{fontspec}
\setmainfont{Latin Modern Roman}
\usepackage{amsmath,amssymb,mathtools,unicode-math}
\setmathfont{Latin Modern Math}
\usepackage{enumitem,needspace,xurl,hyperref}
\setlength{\parindent}{0pt}
\setlength{\parskip}{5pt}
\setlength{\emergencystretch}{3em}
\newcommand{\N}{\mathbb{N}}
\newcommand{\Z}{\mathbb{Z}}
\newcommand{\Q}{\mathbb{Q}}
\newcommand{\R}{\mathbb{R}}
\newcommand{\C}{\mathbb{C}}
\newcommand{\F}{\mathbb{F}}
\newcommand{\A}{\mathbb{A}}
\newcommand{\PP}{\mathbb{P}}
\newcommand{\E}{\mathbb{E}}
\newcommand{\eps}{\varepsilon}
\newcommand{\dd}{\,\mathrm{d}}
\begin{document}
\section*{UnsolvedMath ID: 20001049}
\textbf{Source problem code:} catalogue entry
\subsection{Definitions and mathematical statement}
Together with the preceding local record, this fragment states Conjecture 6.8 under the hypotheses of Seymour--Sullivan Conjecture 6.1; the printed initial ``Conjecture 6.1'' is the continuation of that sentence.

Let $V$ be finite with $n\geq1$ vertices, let $k\geq1$, and let $E_1,\ldots,E_k\subseteq\Omega=\{(u,v)\in V^2:u\ne v\}$. Layer overlap is allowed: an arc has label set $S_{uv}=\{i:(u,v)\in E_i\}$, and $G_i=(V,E_i)$. A rainbow directed cycle admits one distinct chosen label on each of its arcs; no numerical ordering of those labels is imposed. Opposite arcs are allowed in the host, so a rainbow digon is a valid cycle. The conjecture asks that for all nonnegative integers $r_i$,
\[
 \bigl[\delta^+(G_i)\geq r_i\text{ for all }i,\quad
 \textstyle\sum_{i=1}^k r_i\geq n\bigr]
 \ \Longrightarrow\ \text{there is a rainbow directed cycle}.
\]
Equivalently the threshold is $\sum_i\delta^+(G_i)\geq n$. If loops are permitted in another convention, any labeled loop is already a rainbow cycle; the nontrivial case is the loopless one above.

The aid proves the following pairwise special case: for distinct labels $i,j$, if $\delta_i+\delta_j\geq n$, where $\delta_i=\delta^+(G_i)$, then a rainbow digon exists in the loopless case. It reports an exact boundary classification for two layers without a rainbow digon and with $\delta_1+\delta_2=n-1$. Writing $F^R=\{(v,u):(u,v)\in F\}$, those pairs are precisely
\[
 E_1=F,\qquad E_2=\Omega\setminus F^R,
\]
where $F$ is a diregular digraph, meaning every vertex has the same indegree and outdegree $r$. The second layer is then diregular of degree $n-1-r$. Conversely every such $F$ yields a two-label obstruction at total $n-1$. This establishes the two-label threshold; the one-label loopless threshold is impossible since outdegree is at most $n-1$.

For $k\geq2$, the aid's additional incidence bound is
\[
 \sum_{i=1}^k|E_i|\leq k\binom n2
\]
when there is no labeled loop or rainbow digon. The restriction $k\geq2$ matters in this overlapping directed model: with only one label, both directions of every pair may occur without a rainbow digon. Thus the summary's unqualified general-incidence wording is not extended to that exceptional palette. Pairwise obstructions do not prove the conjecture for arbitrary numbers of labels.
\subsection{Short English statement}
The recovered conjecture forces a rainbow directed cycle when minimum layer outdegrees sum to the vertex count. The supplied aid settles the two-label threshold and classifies its sharp boundary, while larger palettes remain outside that special case.
\subsection{Sources}
\begin{itemize}
\item AIM workshop summary on \emph{The Caccetta--H\"aggkvist conjecture}, Conjecture 6.8 using the overlapping-label hypotheses of Seymour--Sullivan Conjecture 6.1: \url{https://aimath.org/WWN/caccetta/caccetta.pdf}.
\item Local registry \texttt{data/unsolvedmath/problems.json}, permanent ID 20001049, exact source code \texttt{catalogue entry}; canonical location \texttt{aim-combinatorics-notes.json notes[173]}. Companion ID 20001048 (\texttt{AIM-COMBINATORICS-0173}, \texttt{notes[172]}) supplies the preceding ``Conjecture 6.8. Under the hypotheses of''; ID 20001036 (\texttt{AIM-COMBINATORICS-0161}) supplies the label-list model.
\item Attached attempt 1, partial result, problem status at run open, rigor proved, low novelty confidence. The two-label classification and incidence bound have no separate complete citation; the aid is machine-generated, not peer reviewed, requiring independent verification before citation.
\end{itemize}
\end{document}
